\chapter{State of the Art}
\label{chap:state-of-the-art}

\section{Introduction}

Securing software has become inseparable from the activity of building it,
and the body of research that surrounds the problem is correspondingly
large. This chapter charts that literature with two aims. The first is to
fix the conceptual vocabulary on which the remainder of the thesis depends,
situating software security within the wider field of computer security and
isolating the two technical problems we address: the detection of
vulnerabilities and their automated repair. The second is to survey the
state of the art for each problem, so that the contribution defended in the
following chapters can be positioned against what already exists. The
chapter then closes by extracting, from this survey, the specific gap that
motivates our work.

\section{Software Security}

The concept of security can be applied to many domains. For the present
chapter, we focus on its application to computing systems and, more
specifically, to the software layer.

Computer security is the discipline that studies and implements mechanisms
to protect computing systems (their data, hardware, and software
components) against unauthorized access, damage, or disruption
\citep{placeholder-sast}. This field can be broken down into several
sub-domains according to the aspect of the computing system being secured,
as illustrated in Figure~\ref{fig:types-computer-security}.

First, there is \texttt{physical security}. Each computing device has a
physical interface through front-panel ports as well as the various
electronic boards. Each of these connections can become a threat if used
maliciously. The most well-known case is the misuse of the USB port or the
deployment of hardware keyloggers \citep{placeholder-detection-survey}.

Another aspect of computer security is \texttt{infrastructure security}.
This branch focuses on the configuration and hardware used to set up an
infrastructure. It includes aspects related to user and access-rights
management as well as workstations. A large part of it focuses on the
network side of the infrastructure, in order to limit each workstation to
its required accesses, while preventing external intrusions and internal
privilege escalations \citep{placeholder-detection-survey}.

Computer security also includes the \texttt{system security} branch that
focuses on operating systems (OS). As the main software component of a
computer, an entire branch of computer security is dedicated to it. The OS
stands out from other software because of its critical role within a machine.
It manages both the hardware aspect of the machine and the users working on
it. Exploiting bugs in an OS can have critical consequences by giving an
attacker full access \citep{placeholder-sast}.

Another branch is \texttt{network security}. Unlike infrastructure security
that deals with the physical setup of the network, network security focuses
on data transmission and its content. Network security can act on the
majority of layers of the OSI model. A well-known example for the general
public is the use of the HTTPS protocol as a replacement for HTTP
\citep{placeholder-sast}.

Computer security also has a more recent branch, halfway between computing
and social sciences: \texttt{social security}. This branch encompasses all
vulnerabilities that use the human factor to penetrate software, networks
or infrastructures. The most common type of attack in this category is
phishing, which relies on the inattention or gullibility of the user to
make them enter confidential information in a corrupted context
\citep{placeholder-detection-survey}.

Finally, we find \texttt{software security}, the branch dedicated to
securing applications, which is associated with software engineering. This
aspect encompasses a large number of practices that involve a variety of
roles, from the developer to the software redistributor. Its level of
criticality has greatly increased in recent years with the democratization
of computing in the population and its intensive use in all areas. The
protection of personal data against theft and resale has become a major
concern. Practices in this field are diverse and can be categorized
according to several criteria, as detailed below \citep{placeholder-detection-survey}.

Figure~\ref{fig:types-computer-security} summarizes the six main
sub-domains of computer security described above. The remainder of this
chapter focuses specifically on the \texttt{software security} branch,
which constitutes the scope of the present work.

\begin{figure}[htbp]
  \centering
  \includegraphics[width=0.9\linewidth]{figures/fig_sec_taxonomy.pdf}
  \caption{Sub-domains of computer security.}
  \label{fig:types-computer-security}
\end{figure}

\subsection{Targets of Software Security}
\label{subsec:targets}

A natural axis for categorizing software security practices is the
\emph{target}: the kind of artefact, runtime and platform on which the
software ultimately executes. Each target carries its own attack surface,
and protective techniques must adapt to that surface rather than be
applied uniformly. Modern software production has multiplied these
targets, and any taxonomy that aspires to be comprehensive has to keep
pace with this diversity \citep{placeholder-detection-survey}. Rather
than presenting the targets as a long flat list, Table~\ref{tab:sec-targets}
first gives a compact map of the main execution contexts and the attack
surface that dominates each of them.

\begin{table}[htbp]
\centering
\small
\renewcommand{\arraystretch}{1.18}
\begin{tabularx}{\linewidth}{@{}p{0.24\linewidth}p{0.31\linewidth}X@{}}
\toprule
\textbf{Target} & \textbf{Typical context} & \textbf{Dominant attack surface} \\
\midrule
Desktop applications & Personal computers and general-purpose operating systems & Memory safety, local privilege, local parsers, supply chain \\
Server-side / back-end & Physical or virtualised servers exposed through APIs & Injection, broken authentication, access control, SSRF \\
Web applications & Browser-based clients and HTTP services & XSS, CSRF, clickjacking, DOM attacks, parameter tampering \\
Mobile applications & Android, iOS and cross-platform runtimes & Permission misuse, insecure storage, inter-component communication, repackaging \\
Embedded / IoT & Constrained devices and firmware-oriented deployments & Memory corruption, weak update mechanisms, limited mitigations \\
Cloud-native / serverless & Containers, functions, micro-services and orchestrators & Misconfiguration, container images, secrets, infrastructure-as-code \\
OS / system-level & Kernels, hypervisors, drivers and runtime engines & Privileged code execution and subversion of higher-level protections \\
Smart contracts & Blockchain virtual machines such as the EVM & Reentrancy, integer overflows, state manipulation, bridge flaws \\
Third-party dependencies & Open-source libraries and transitive package graphs & Propagated vulnerabilities, vulnerable components, reachability \\
Machine-learning systems & Models, prompts and data-processing pipelines & Model backdoors, prompt injection, extraction, unsafe artefact loading \\
\bottomrule
\end{tabularx}
\caption{Software-security targets, contexts and dominant attack surfaces.}
\label{tab:sec-targets}
\end{table}

The table can be read as four broad families. \emph{Client-facing}
targets, such as desktop, web and mobile applications, are shaped by
user-controlled inputs, local permissions and the trust boundary between
the user interface and the service it contacts. \emph{Service-side}
targets, including back-end, cloud-native and serverless workloads, expose
their attack surface through APIs, configuration, authentication,
authorization and secrets management. \emph{Low-level and specialized}
targets, such as embedded software, operating systems and smart contracts,
are harder to secure because a single flaw can bypass many of the
protections available at higher layers. Finally, \emph{inherited and
emerging} targets, namely third-party dependencies and machine-learning
systems, show that modern applications also inherit risks from components,
models and data pipelines that are not always developed by the application
team itself \citep{placeholder-sast,placeholder-sca}.

The taxonomy presented in the rest of this thesis is deliberately kept
platform-agnostic. The detection and remediation techniques examined later
can therefore be applied, with minor adaptations, across most of these
targets whenever the relevant source code or design artefact is available.

\subsection{Stages of Security Protocol Application}
\label{subsec:stages}

Another criterion for categorizing the software security domain is the
\emph{timing} of its implementation. The development and deployment cycle
of an application is complex, and a piece of software goes through several
states in which it is possible to intervene to secure it
\citep{placeholder-sast}:

\texttt{Source state:} this is the state in which the application is
during the development phases. In this state, the code is easily
understandable by a human. During these phases, software security takes the
form of good practices to be applied by developers, recommendations
regarding vulnerabilities in the code, or the use of external tools that
modify the code \citep{placeholder-sast}.

\texttt{Packaged state:} a packaged application is one that is ready to
be deployed. It is not necessarily in the form of a single package: it can
be the result of one or more optimizers, compilers, transpilers and
packagers used to make it usable in production. The actual form of the
application depends heavily on the technologies used and the target
platform. In this state, the code is generally hardly readable by a human
and is rather targeted at a physical or virtual machine. Software security
is often linked here to approaches implemented by external tools that
modify the compiled code \citep{placeholder-sast}.

\texttt{Executed state:} this is the state in which the software is when
it is in use on a machine. Software security then takes the form of internal
security mechanisms within the application, or of an external application
observing and controlling the execution \citep{placeholder-detection-survey}.

\subsection{Actions of Software Security}

A last criterion for categorizing software security is the type of action
performed to defend the application. These actions constitute the
operational layer of the taxonomy and apply, in varying forms, across the
targets introduced in Section~\ref{subsec:targets} and at each of the
stages described in Section~\ref{subsec:stages}. They can be carried out
either automatically or manually \citep{placeholder-sast}.

\begin{itemize}
  \item \texttt{Malware identification:} the detection of intentionally
        malicious software. It aims to protect applications by detecting
        and processing external threats \citep{placeholder-detection-survey}.
  \item \texttt{Vulnerability detection:} it can be done from two angles.
        The first is the \emph{internal} angle, which aims at detecting
        vulnerabilities left in the application's code. The second is the
        \emph{external} angle, which aims at detecting vulnerabilities in
        the execution environment, that is, weaknesses that can present a
        risk to the system's integrity in the face of an attack. The
        external angle is by nature limited to security applied at runtime
        since an environment is required to detect vulnerabilities in it.
        It particularly applies to client applications, for which the
        environment is unknown and not controlled \citep{placeholder-detection-survey}.
  \item \texttt{Remediation:} this approach aims at modifying the code of
        an application in order to remove its vulnerabilities. It relies
        on the previously mentioned vulnerability detection. Two angles
        are possible: \emph{corrective} remediation, where an existing
        vulnerability has to be removed from the code; and
        \emph{preventive} remediation, which consists in hardening the
        application beforehand to make exploitation more difficult,
        even when no specific flaw has been pinpointed. The present work
        focuses on the corrective angle \citep{placeholder-apr-survey}.
  \item \texttt{Runtime Application Self-Protection (RASP):} an action
        specific to the executed state, whose goal is to protect an
        application during its execution. It can either embed additional
        security functions to be executed alongside the application, or
        rely on an external application monitoring the application to be
        secured during its execution \citep{placeholder-iast}.
  \item \texttt{Testing:} this approach uses test suites that perform
        security checks. They take the same form as unit or integration
        tests, and accordingly share their advantages and drawbacks. They
        can be performed in white-box, gray-box or black-box settings,
        depending on the level of knowledge of the internal behavior of
        the application \citep{placeholder-sast}.
\end{itemize}

The two actions of greatest relevance to the present work are vulnerability
detection and remediation. The techniques underpinning each of them are
examined in detail in the following two subsections.

\section{Vulnerability Detection Techniques}
\label{subsec:detection-techniques}

Before surveying detection techniques, it is necessary to establish a
working definition. In the context of this work, a \emph{vulnerability} is
a weakness in a software system, arising from a design flaw, an
implementation error, or a misconfiguration, that can be exploited by a
malicious actor to compromise the confidentiality, integrity, or
availability of the system or its data \citep{placeholder-detection-survey}.

Detecting a vulnerability before an attacker exploits it is the prerequisite
of any remediation effort. The literature has produced a rich and rather
fragmented body of techniques, which we group below according to the
principle that drives the analysis rather than the tooling that implements
it. The boundaries between families are porous: many modern tools combine
several paradigms, and recent learning-based detectors blur the classical
static/dynamic distinction altogether \citep{placeholder-detection-survey}. The
principal families are presented in turn below. These techniques are presented in the following:

\texttt{Static Application Security Testing (SAST):}
Static analysis inspects the source code, the bytecode or the binary
without executing it. The simplest form relies on syntactic
pattern matching, looking for code constructs known to be dangerous,
for instance a call to \texttt{strcpy} in C or a SQL query built by string
concatenation. More sophisticated engines abstract the program into
intermediate representations such as control-flow graphs, data-flow graphs
or call graphs, then run analyses like taint tracking, symbolic execution
or abstract interpretation \citep{placeholder-sast}. The strength of SAST
lies in its early-stage applicability, well before any test campaign;
its weakness is a non-negligible rate of false positives, which has
historically discouraged adoption in some industrial contexts. Widely used
representatives include SonarQube, Coverity, Checkmarx, Fortify and the
academically influential FlawFinder and Infer projects.

\texttt{Dynamic Application Security Testing (DAST):}
Dynamic techniques execute the program and observe its behaviour. Unlike
SAST, they require either a deployable build or a running service. The
canonical instance is \emph{fuzzing}, where malformed or random inputs are
fed to the program in order to trigger crashes, unexpected exceptions or
memory corruptions. Coverage-guided fuzzers such as AFL, libFuzzer and
their many descendants \citep{placeholder-fuzzing} have become a standard
in low-level software, while web-oriented scanners (OWASP ZAP, Burp Suite)
target HTTP-level flaws like SQL injection or cross-site scripting.
Compared with SAST, DAST suffers far fewer false positives, but it can
only flag what its execution actually reaches: code paths that the test
harness does not visit remain invisible.

\texttt{Interactive Application Security Testing (IAST):}
IAST combines instrumentation with runtime observation: agents inserted in
the application correlate the inputs received by the program with the
sensitive operations they end up reaching. The result is a hybrid view
that retains the precision of dynamic analysis while preserving some of
the coverage advantages of static reasoning \citep{placeholder-iast}. The
trade-off is the deployment cost, IAST agents must be embedded in the
runtime, and a stronger dependency on representative workloads.

\texttt{Software Composition Analysis (SCA):}
Modern systems are assembled rather than written: a typical project pulls
hundreds of third-party dependencies through package managers such as npm,
Maven or PyPI. SCA tools enumerate this dependency graph and cross-check
each component against vulnerability databases, the most widely used being
the NVD and the GitHub Advisory Database \citep{placeholder-sca}. The
challenge here is less about detection logic than about the quality of
the underlying advisories and the difficulty of distinguishing
\emph{vulnerable code present in the binary} from \emph{vulnerable code
actually reachable from the application}.

\texttt{Learning-based detection:}
A more recent line of work applies machine learning to the problem,
treating vulnerability detection either as a classification task on code
fragments or as an anomaly-detection task on traces. Early approaches
used hand-crafted features feeding classical classifiers (SVM, random
forests). A second wave moved to deep learning, with recurrent and
convolutional networks operating on tokenised code, then to graph neural
networks operating on AST or program-dependence graphs, Devign and
Reveal being often-cited representatives \citep{placeholder-gnn}. The
current frontier is the use of large language models, either as zero-shot
detectors prompted with a vulnerability description, or as fine-tuned
classifiers trained on labelled corpora such as Big-Vul or D2A
\citep{placeholder-llm-detection}. These models extend coverage to
vulnerabilities for which no syntactic signature exists, but they
inherit well-known weaknesses: opacity, dataset biases, and a tendency
to memorise rather than generalise.

\texttt{Hybrid approaches:}
A non-trivial fraction of recent works deliberately combine several of
the families above. A common pattern is to use SAST for a first coarse
pass, then DAST or symbolic execution to confirm the candidates flagged
statically, in order to filter false positives without losing recall.
Another pattern feeds the output of a static analyser as input features
to a learning model, letting the network re-rank or prune the warnings
\citep{placeholder-hybrid} .

\texttt{Classification by input artefact:}
The families above are best read along a second, orthogonal axis: the
artefact each technique consumes as input. This axis is not merely
taxonomic; it directly determines \emph{when} in the development
lifecycle detection can take place, and therefore how early a discovered
flaw can be corrected.

The dominant input is the program \emph{source code}. SAST operates on
source, bytecode or binary without execution, making it applicable as
soon as the developer writes the first lines of code. Learning-based
detectors trained on source corpora similarly require nothing more than
the code itself, and can therefore be integrated into a continuous
integration pipeline at the earliest stages of development. Both
families thus belong to the \emph{pre-deployment} tier: they surface
weaknesses before the application is built or shipped, when the cost of
correction is lowest.

A second category of techniques takes a \emph{running program} as
input. DAST requires a deployable build, and IAST requires an
instrumented runtime; both therefore intervene later, once the software
has reached a testable or staging state. SCA operates on the
\emph{dependency graph} rather than on the program itself, and can be
run at build time, making it an early-stage technique in practice
despite its different input form.

Hybrid approaches, by combining artefacts from more than one of these
categories, can span multiple stages but also accumulate the deployment
constraints of each component they include.

Table~\ref{tab:detection-families} summarizes the principal families
along this input-artefact axis together with the development stage at
which each can intervene and its main trade-off.

\begin{table}[htbp]
\centering
\small
\renewcommand{\arraystretch}{1.15}
\begin{tabular}{llll}
\hline
\textbf{Family} & \textbf{Input artefact} & \textbf{Stage} & \textbf{Main trade-off} \\
\hline
SAST           & Source / bytecode / binary  & Development    & Early use, higher false positives \\
Learning-based & Source code or traces       & Development    & Wide coverage, opacity and bias \\
SCA            & Dependency graph            & Build time     & Advisory quality, reachability \\
DAST           & Running program             & Testing        & Few false positives, limited coverage \\
IAST           & Instrumented runtime        & Testing        & Precise, costly deployment \\
Hybrid         & Combined artefacts          & Varies         & Higher recall, integration cost \\
\hline
\end{tabular}
\caption{Detection families classified by input artefact and development stage.}
\label{tab:detection-families}
\end{table}

Figure~\ref{fig:detection-taxonomy} summarizes the main families of
vulnerability detection techniques described above. It illustrates how
the learning-based sub-family has progressively diversified from
classical machine learning towards deep learning, graph neural networks
and large language models, while remaining distinct from hybrid
approaches that combine learning with non-learning paradigms.

\begin{figure}[htbp]
  \centering
  \includegraphics[width=0.92\linewidth]{figures/fig_detection_taxonomy.pdf}
  \caption{Families of vulnerability detection techniques.}
  \label{fig:detection-taxonomy}
\end{figure}

Having established the landscape of detection techniques, the following
section turns to the complementary problem of repairing the
vulnerabilities once detected, which is the domain of Automated Program
Repair.


\section{Automated Program Repair}
\label{subsec:apr}
\textcolor{green}{Is this a subsection of the section named "Vulnerability Detection Techniques ?"}
Automated Program Repair (APR) gathers the techniques whose goal is to
correct defects in a program automatically, without altering its expected
behaviour \citep{placeholder-apr-survey}. Historically the field grew out
of the bug-fixing literature, and the term \emph{defect} was nearly
synonymous with \emph{functional bug}. The notion has since been
broadened: a defect can equally denote a code-quality smell, a
performance regression or, as in our case, a security vulnerability
\citep{placeholder-apr-vuln} . Adopting this wider definition is what
makes APR directly relevant to the second half of our research question,
since remediating a vulnerability and repairing a bug share the same
abstract structure: locating a faulty fragment, producing a candidate
modification, and convincing oneself that the modification is sound.

Most APR pipelines, regardless of the underlying technology, decompose
the problem into the same three stages \citep{monperrus2018apr-survey}:
\emph{fault localization} narrows the search down to a code region;
\emph{patch generation} proposes one or more candidate modifications for
that region; \emph{patch validation} decides which candidates, if any,
actually solve the problem without breaking anything else. The three
stages are not strictly sequential in practice. Learning-based pipelines,
in particular, often blur the boundary between the first two by emitting
a localized patch in a single inference step, and the validation feedback
can be looped back into either of the upstream stages.

\begin{figure}[htbp]
  \centering
  \includegraphics[width=0.92\linewidth]{figures/fig_apr_pipeline.pdf}
  \caption{Canonical stages of an Automated Program Repair pipeline.}
  \label{fig:apr-pipeline}
\end{figure}

\begin{enumerate}

\item \texttt{Fault Localization (FL):} is the entry point of the pipeline. Its task is
to point at the fragment of code that is responsible for the observed
defect, with a granularity that ranges, depending on the technique, from
an entire file down to a single statement \citep{wong2016fl-survey}.
Granularity matters: a coarse output forces the patch generator to
explore a much larger search space, while an overly fine one is
vulnerable to off-by-one errors that send the rest of the pipeline
chasing the wrong line.

Two failure modes deserve to be highlighted. A false positive, a
location flagged as faulty when it is not, typically translates into
wasted computation and, in the worst case, into a regression introduced
by a useless edit. A false negative is more damaging: the real defect
survives the run untouched, and the operator may end up with a
false sense of security. Some learning-based methods avoid the question
altogether by entangling localization and patch generation into a single
end-to-end model; this simplifies the engineering at the cost of
interpretability.

\item \texttt{Patch Generation(PG): } sits at the heart of APR. Given a localized
defect, the generator produces one or several candidate modifications.
The space of plausible patches is enormous, a single line can be
rewritten in countless syntactically valid ways, and exploring it
naively does not scale. The art of PG is therefore as much about
\emph{pruning} as it is about \emph{producing}: heuristics, learned
priors, repair templates and constraint solvers all serve, in their own
way, to keep the candidate set tractable \citep{placeholder-pg} . Time
budget is a hard requirement in any realistic deployment, and a method
that finds the right patch in principle but only after several hours of
computation is of limited practical value.

\item \texttt{Patch Validation (PV): } closes the loop. Each candidate produced upstream
must be checked on two counts: that it actually removes the targeted
defect, and that it does not introduce a new one elsewhere in the
program. The dominant approach in the literature is test-suite-based
validation, which inherits both the strengths and the weaknesses of
functional testing: it is cheap and reproducible when the tests exist,
but its verdict is only as exhaustive as the suite itself. The
well-known phenomenon of \emph{patch overfitting}, patches that pass
all available tests yet fail to capture the real fix, is a direct
symptom of this limitation \citep{qi2015overfitting} . Formal techniques
offer stronger guarantees, but their cost in expertise and proof effort
has so far prevented their generalization.

    
\end{enumerate}



\subsection{Systematic Mapping Studies}
\label{subsec:sms}

The two domains presented above, detection and repair, are highly
active research fields. Hundreds of primary studies are published every
year, employing heterogeneous vocabulary, evaluating on incompatible
benchmarks and reporting partial results. Before any new contribution
can be designed, the literature itself has to be charted. This is the
role of \emph{secondary studies}, and more specifically of
\emph{Systematic Mapping Studies (SMS)} \citep{petersen2015sms} .

A SMS aims at providing a structured overview of a research area by
classifying its primary studies along a set of well-defined dimensions,
typically including the type of contribution, the underlying technique,
the addressed problem and the empirical setting. The output is less an
in-depth synthesis, which is the province of a \emph{Systematic
Literature Review (SLR)}, than a map: a categorical view of where the
field is dense, where it is thin, and which combinations of dimensions
have not yet been explored \citep{kitchenham2009slr} . The distinction
matters in our context, because the question we ultimately ask,
\emph{what techniques exist, for which vulnerabilities, and how do
detection and remediation interact?}, is fundamentally a mapping
question, not a meta-analytic one.

A SMS unfolds in five stages, broadly stable across the methodological
literature \citep{petersen2008sms-original}. \emph{Planning} fixes the
research questions, the inclusion and exclusion criteria, and the search
strategy. \emph{Search} executes that strategy on a chosen set of
digital libraries, possibly extended by snowballing on the references
and citations of seed papers. \emph{Screening} reduces the raw search
output to the actually relevant subset, usually through a sequence of
title-and-abstract filtering followed by full-text inspection.
\emph{Extraction} pulls structured data out of each retained paper
according to the predefined classification scheme. \emph{Synthesis}
finally aggregates this data into the maps, tables and narrative
sections that constitute the report.

\begin{figure}[htbp]
  \centering
  \includegraphics[width=0.86\linewidth]{figures/fig_sms_process.pdf}
  \caption{Canonical stages of a Systematic Mapping Study.}
  \label{fig:sms-process}
\end{figure}

Conducted manually, this process is notoriously labour-intensive. A
single SMS routinely involves screening several thousand candidate
papers and extracting dozens of fields per retained study, with two
reviewers working in parallel to limit subjective bias. Three
difficulties recur in the literature \citep{placeholder-sms-pain} .
\emph{Scalability} first: the volume of publications in software
security has grown super-linearly over the last decade, to the point
that an exhaustive screening is no longer feasible without tooling.
\emph{Reproducibility} next: small variations in query strings, in the
timing of the search or in the interpretation of the inclusion criteria
can yield substantially different corpora, and the standard practice of
reporting Cohen's $\kappa$ between reviewers, while useful, does not
fully address the issue. \emph{Currency}, finally: by the time a SMS is
published, a non-trivial fraction of its corpus is already a year or
two old. These pain points are precisely what motivates the recent
interest in (semi-)automated SMS pipelines, which we discuss in
Section~\ref{subsubsec:sms-auto}.

\subsection{Towards Automated Mapping Studies}
\label{subsubsec:sms-auto}

The pain points just listed have triggered a wave of tools that try to
inject automation into the SMS workflow without giving up its
methodological guarantees. Earlier systems such as ASReview, Rayyan,
EPPI-Reviewer or RobotReviewer mostly assist a single phase, typically
abstract screening, deduplication or evidence management, while the
reviewer still performs the decisive methodological work. More recent
approaches delegate individual screening calls to a large language
model, but they tend to bind the prompt, the inclusion logic and the
classification schema to one specific review protocol. They therefore
reduce effort locally without providing a reusable, end-to-end SMS
pipeline \cite{placeholder-sms-auto}.

SMS-Auto is a human-in-the-loop LLM pipeline that
automates the operational stages of an SMS while preserving researcher
control at each point where methodological judgement is required. The
following subsections describe the problem addressed by SMS-Auto, its
layered architecture, its prompt synthesis mechanism, its screening
protocol, the downstream extraction and categorization stages, and the
empirical evidence reported for its reliability. Figure~\ref{fig:smsauto-pipeline}
gives an overview of the automated stages and the mandatory human
checkpoints that punctuate them.

\begin{figure}[htbp]
\centering
\includegraphics[width=0.96\linewidth]{figures/fig_smsauto_pipeline.pdf}
\caption{End-to-end SMS-Auto pipeline.}
\label{fig:smsauto-pipeline}
\end{figure}

\subsubsection{Rationale for Automation}

The starting point of SMS-Auto is not that SMS methodology is weak, but
that it is expensive precisely because it is rigorous. A protocol must
define the study objective, research questions, Boolean search strings,
inclusion and exclusion criteria, extraction fields and classification
scheme before any paper is processed. Once the search is executed, a
single SMS can easily produce several thousand candidate records. The
paper reports that screening alone may require several reviewer hours
per small batch of papers, before deduplication, snowballing,
full-text extraction and synthesis are even considered
\cite{nouali2026smsauto}. For fast-moving areas such as software
security, this delay creates a practical paradox: by the time a manual
mapping study is ready for publication, part of the corpus may already
be outdated.

The automation target is therefore carefully delimited. SMS-Auto does
not attempt to replace the methodological responsibility of the
researcher. Instead, it automates the repetitive and computationally
amenable operations: collecting records from scholarly APIs, detecting
duplicates, filtering clearly irrelevant papers, running independent
screening decisions, expanding the corpus through snowballing,
extracting structured evidence from full texts, and grouping papers
according to user-defined dimensions. Human judgement remains attached
to protocol design, checkpoint validation, conflict resolution and final
interpretation. This division of labour is essential, because an SMS is
not merely an information-retrieval task; it is a reproducible research
method.

\subsubsection{Layered Architecture}

SMS-Auto is organised into five layers whose separation is the main
reason the system can be reused across topics. The first layer, User
Interaction, treats the researcher as an actor inside the pipeline
rather than as a passive consumer of exported files. At initialization,
the researcher supplies the protocol object: objective, research
questions, search strings, data sources, inclusion criteria and
exclusion criteria. During execution, the same interface exposes five
mandatory checkpoints: after collection, automated exclusion,
dual-model screening, snowballing and extraction. At each checkpoint,
the researcher can inspect the intermediate corpus, override decisions,
or restart a stage.

The second layer, Workflow Orchestration, maintains the state of the
SMS lifecycle. It tracks whether a paper has been retrieved,
deduplicated, excluded, screened, manually reviewed, extracted or
categorized. It also enforces the checkpoint gates. A downstream stage
cannot proceed simply because the previous computation finished; it
can proceed only after the corresponding human checkpoint has been
validated. When a stage is restarted, the orchestrator invalidates the
dependent downstream states so that the exported artefacts remain
consistent with the latest protocol and paper set.

The third layer, Protocol and Prompt Management, contains the
Protocol-to-Prompt engine. This is the architectural centre of the
system, because it decouples study logic from model execution logic.
The fourth layer groups the Core Processing Services: source retrieval,
deduplication, automatic exclusion, LLM screening, conflict resolution,
snowballing, extraction, categorization and export. Each service has a
narrow responsibility, which makes it easier to validate and to replace
independently. The fifth layer contains External Services: metadata
APIs such as OpenAlex and CrossRef, citation graph APIs such as
Semantic Scholar, LLM providers such as Cerebras and Groq, and PDF
retrieval sources. The implementation includes retry logic, exponential
backoff and API-key rotation to keep long runs robust under rate limits
\cite{nouali2026smsauto}.

\begin{figure}[htbp]
\centering
\includegraphics[width=0.86\linewidth]{figures/fig_smsauto_arch.pdf}
\caption{Five-layer architecture of SMS-Auto.}
\label{fig:smsauto-arch}
\end{figure}

\subsubsection{Protocol-to-Prompt Generalization}

The distinctive contribution of SMS-Auto lies in the way prompts are
produced. A fragile implementation would write one prompt for a
vulnerability-detection SMS, another for a remediation SMS, and so on.
Such a design would automate a single study but not the SMS method. In
SMS-Auto, the prompts are synthesized deterministically from the
protocol object supplied by the researcher. The objective and research
questions are rendered into a scope description for the exclusion
stage. The inclusion and exclusion criteria are converted into numbered
decision rules for screening. The extraction fields become a typed JSON
contract for evidence extraction, and the categorization dimensions
become grouping instructions for synthesis.

Internally, the prompt builder follows a three-step chain. It first
loads a stage-specific skeleton defining the model role, task and
expected output format. It then injects the relevant protocol fields
with type checks and length guards. Finally, it appends consistency
constraints and a structured output schema. Because this procedure is
deterministic, the same protocol and the same paper always generate
the same prompt. This point matters methodologically: reproducibility
in an LLM-assisted SMS depends not only on preserving the final corpus,
but also on preserving the exact instructions that produced each
intermediate decision.

\subsubsection{Screening and Conflict Resolution}

The first model-driven stage is automatic exclusion. It is deliberately
designed as a coarse, high-recall filter operating on titles and
abstracts. Its task is to remove records that are clearly out of scope,
not to make the final inclusion decision. The prompt therefore uses a
fail-open policy: when the model is uncertain, it must keep the paper
for downstream screening. This asymmetry favours recall over precision,
which is appropriate at the beginning of an SMS because a paper kept by
mistake can still be removed later, whereas a relevant paper removed
too early is lost.

The second model-driven stage is dual-LLM screening. Each candidate is
evaluated independently by two LLM instances hosted on different
inference providers. In the reported implementation, both instances use
the Llama 3.3 70B family, with one served by Cerebras and the other by
Groq. This choice is intentional: it keeps the base model stable while
exposing the decision to different provider infrastructures. The
screening prompt includes the reviewer role, the full inclusion and
exclusion criteria, consistency rules, selected paper content and a
fixed JSON schema containing the decision, applied criteria,
justification and confidence.

Conflict resolution is tri-state. If both models agree and both
confidence scores are at least $0.7$, the consensus decision is
accepted automatically. If they agree but one confidence score is below
$0.7$, the consensus is retained but the paper is flagged for manual
review. If the models disagree, the paper is assigned
\texttt{NEEDS\_MANUAL\_REVIEW} regardless of confidence. This protocol
treats disagreement as useful evidence rather than noise. The reviewer
sees both justifications and arbitrates only those cases where the
automation has surfaced genuine ambiguity.

\begin{figure}[htbp]
\centering
\includegraphics[width=0.9\linewidth]{figures/fig_smsauto_screening.pdf}
\caption{Dual-LLM screening and tri-state conflict resolution.}
\label{fig:smsauto-screening}
\end{figure}

\subsubsection{Snowballing, Extraction and Categorization}

After the screening checkpoint, SMS-Auto expands the validated corpus
through backward and forward snowballing. Backward snowballing follows
the references of included papers; forward snowballing follows later
papers that cite them. The paper reports the use of multiple scholarly
services and merged results rather than a single fallback chain, since
citation coverage differs substantially between providers. This stage
is especially important in software engineering, where keyword search
alone can miss relevant papers because terminology changes quickly or
because different communities name the same technique differently.

The extraction phase starts once the snowballed corpus has passed its
checkpoint. For every included paper, the system retrieves an
accessible PDF, extracts the text with pdfplumber and falls back to
PyPDF2 when needed. Long texts are truncated at sentence boundaries to
fit the LLM context window. The extraction prompt asks for a structured
JSON object containing metadata, synthesis paragraphs, technical
attributes, datasets, reported results, limitations, contributions and
research gaps. Negative instructions are explicit: the model must not
infer results that are absent from the paper, and it must return
\texttt{null} when a field cannot be answered from the text.

Categorization is the final analytical layer. The researcher can define
several independent dimensions, for example method family, application
domain, evaluation strategy, target language or vulnerability class.
For each dimension, the system asks the model to assign papers to
clusters and to produce a short rationale for the grouping. It then
aggregates the extracted evidence inside each cluster, which gives the
researcher material that can be reused in tables and narrative
synthesis. The validation study also identifies a limitation of this
stage: open-vocabulary LLM categorization tends to produce labels that
are too specific. For this reason, taxonomy-constrained prompts are
needed when the purpose is to map papers onto a predefined scheme
rather than discover categories inductively.

\subsubsection{Empirical Validation}

SMS-Auto was validated through five experiments covering the main
LLM-driven stages of the pipeline. The reported results are summarized
in Table~\ref{tab:sms-auto-validation}. The exclusion experiment shows
why prompt design matters: an overly strict first prompt achieved poor
recall, while a broad second prompt kept too many irrelevant papers.
The balanced production prompt reached recall $0.915$, $F_1=0.714$ and
Cohen's $\kappa=0.50$, which is acceptable for a first-pass high-recall
filter.

\begin{table}[htbp]
\centering
\small
\begin{tabular}{lll}
\hline
Experiment & Stage & Main result \\
\hline
V3 & Automatic exclusion & Recall $0.915$, $F_1=0.714$ \\
V4 & Dual screening & Flip rate $0\%$, Fleiss' $\kappa=1.000$ \\
V5 & Snowballing & Backward recall $0.705$, forward recall $0.894$ \\
V6 & Extraction & Hallucination $0\%$, BERTScore $0.876$ \\
V7 & Categorization & LLM stability $\kappa=0.894$ \\
\hline
\end{tabular}
\caption{SMS-Auto validation results}
\label{tab:sms-auto-validation}
\end{table}

The screening experiment is the strongest evidence for operational
stability. Across three independent runs on the same one hundred-paper
sample, no binary decision flipped, Fleiss' $\kappa$ reached $1.000$,
and confidence scores remained stable with an intra-class correlation
of $0.954$. The extraction experiment is equally important for trust:
metadata fields such as year, venue and DOI reached full accuracy after
normalization, title accuracy reached $90\%$, and numeric results had a
hallucination rate of $0\%$. The weakest result is not a failure of
screening but a categorization alignment issue: the model was stable
with itself, yet its labels were more fine-grained than the expert
taxonomy, leading to low LLM-to-expert agreement until label
constraints were added.

The practical effect is a reduction of the manual burden, not the
removal of the reviewer. In the companion cybersecurity mapping work,
the detection SMS processed $14{,}827$ raw records and produced a
validated corpus of $99$ papers in approximately three calendar weeks
with two researchers. The authors estimate that deduplication and
screening alone would otherwise have required an additional four to six
weeks \cite{nouali2026smsauto}. This is the contribution that matters
for the present thesis: SMS-Auto makes a systematic map feasible within
the time constraints of a research project, while keeping the protocol,
checkpoints and final interpretations under human control.




\section{Related Work}
\label{sec:related-work}

With the conceptual scope now fixed, we turn to the works that share its
objective. The literature on automated detection and remediation of
software vulnerabilities is large but unevenly distributed: a handful of
techniques attract a disproportionate share of the publications, while
several sub-areas remain thinly populated. The aim of this section is
not to be exhaustive, the two systematic mapping studies introduced
later in this work carry that responsibility, but to give a structured
overview of the sub-fields directly adjacent to our contribution.

We organise the discussion along the two halves of our research scope.
Section~\ref{subsec:rw-detection} surveys works that aim at
\emph{detecting} vulnerabilities in source or binary code, grouped by
analysis paradigm rather than by tool name. Section~\ref{subsec:rw-apr}
turns to \emph{remediation}, where Automated Program Repair has been
the dominant framing for the last decade.
Section~\ref{subsec:integrated-rw} covers the few works that try to
close the loop between the two, and
Section~\ref{subsec:secondary-rw} positions our effort against the
existing secondary studies. The empirical counts we cite below come
from the systematic mapping reported later in this thesis, and are
introduced here only to give a sense of how the literature is
distributed across paradigms.

\subsection{Vulnerability Detection}
\label{subsec:rw-detection}

Section~\ref{subsec:detection-techniques} introduced the principal
families of detection techniques from a conceptual and taxonomic
perspective, describing what each paradigm does and what properties it
offers. The present section complements that view by surveying the
\emph{specific works} that instantiate those paradigms: concrete tools,
datasets and empirical studies published in the literature. The
distinction is therefore between a classification of techniques and a
review of contributions; the former structures the landscape, the
latter populates it with evidence.

A first observation when surveying the detection literature is that
the boundaries between paradigms are blurrier than introductory
expositions suggest. Many recent tools combine several techniques, and
the labels we use below should be read as identifying the
\emph{dominant} mechanism rather than an exclusive one. At a high
level, the surveyed corpus splits into three blocs: classical
program-level analyses, AI- and learning-based detectors, and hybrid
combinations of the two.

\begin{enumerate}


\item \texttt{Classical Program-Level Analyses}
\label{subsubsec:rw-classical}

\begin{itemize}

  \item \textit{Static analysis.} Static analysis remains the most
        populated classical family. Recent contributions cover both
        general-purpose and domain-specific settings: D2A builds a
        labelled dataset of real-world bugs from differential analysis
        of OSS commits \citep{placeholder-d2a}; Eth2Vec, SmartAxe and
        SmartState target smart contracts on the Ethereum Virtual Machine
        \citep{placeholder-eth2vec,placeholder-smartaxe}; argXtract and a
        wave of comparative evaluations of mainstream SAST tools
        (SonarQube, Coverity, CodeQL) study how well static analysers
        transfer between projects and languages
        \citep{placeholder-sast-eval}. The recurring weaknesses reported
        by these works, namely high false-positive rates, sensitivity to
        coding conventions, and difficulty in modelling inter-procedural
        and inter-component flows, are the same that motivated the move
        towards learning-based detectors.

  \item \textit{Taint analysis.} Several studies refine the
        long-standing taint-tracking framework. TaintBench provides a
        reproducible benchmark for Android
        \citep{placeholder-taintbench}; Taint Rabbit and the
        ``bimodal taint'' family revisit precision-vs-coverage
        trade-offs at binary level \citep{placeholder-taint-rabbit};
        more recent work pairs taint propagation with LLM-driven sink
        detection on web back-ends, with reported $F_1$ scores between
        $0.76$ and $0.97$ on internal datasets
        \citep{placeholder-llm-taint}. Taint analysis stays attractive
        when the threat model is dominated by data-flow vulnerabilities
        (injection, information leaks), but its dependence on a precise
        model of the runtime keeps it expensive to deploy on
        heterogeneous codebases.

  \item \textit{Program and binary analysis.} A further set of studies
        groups around program analyses that are neither pure SAST nor
        pure taint, labelled either ``program analysis'' (e.g.\
        slicing-based clone detection, LEOPARD's program-metric approach
        to vulnerable-code identification \citep{placeholder-leopard}) or
        ``binary analysis'' (memory-safety enforcement \`{a} la
        EffectiveSan \citep{placeholder-effectivesan},
        opaque-predicate hardening, buffer-overflow refinement). They
        share a code-centric, semantics-aware view but differ in the
        artefact they reason about, source vs.\ compiled binaries,
        which has consequences on what can be repaired afterwards.

  \item \textit{Web application security testing.} A distinct cluster
        of studies addresses vulnerabilities specific to the web stack,
        ranging from XSS and access-control flaws in cloud add-ons to
        parameter tampering and server-side request forgery. Recent
        entries in this cluster (e.g.\ Artemis on SSRF
        \citep{placeholder-artemis}, MalCodeAI on multi-class detection
        and remediation \citep{placeholder-malcodeai}) increasingly
        bridge into the learning-based and hybrid families discussed
        below, suggesting that classical web-AST tooling alone may have
        reached a plateau.

\end{itemize}

\item \texttt{AI and Learning-Based Detection}
\label{subsubsec:rw-ai}

It is important here to insist on a hierarchy that the literature
sometimes blurs. \emph{Machine learning}, \emph{deep learning},
\emph{graph neural networks} and \emph{large language models} are not
four parallel detection paradigms; they are successive layers of the
same family, the AI- or learning-based family. We therefore treat them
as sub-categories of a single block, and reserve the term
\emph{hybrid} for approaches that genuinely combine learning with a
non-learning paradigm.

Within this family, deep learning dominates the corpus. A large body
of work falls under this label, and the spread of techniques mirrors a
decade of progress in the broader ML field. Early systems trained
recurrent or convolutional networks on tokenised code slices labelled
with CWE identifiers, with VulDeePecker and SySeVR as canonical
examples \citep{placeholder-vuldeepecker,placeholder-sysevr}. A second
wave moved to graph neural networks operating on AST, control-flow or
program-dependence graphs, Devign \citep{placeholder-devign}, Reveal
\citep{placeholder-reveal}, LIVABLE for long-tailed CWE classification
\citep{placeholder-livable}, exploiting longer-range dependencies that
flat-token models systematically miss. The third and most recent wave
relies on transformer-based language models: VulBERTa, XGV-BERT,
EnStack and DSHGT either fine-tune encoders such as CodeBERT on
labelled corpora (Big-Vul, D2A) or stack several LLMs to improve
recall on under-represented classes
\citep{placeholder-vulberta,placeholder-xgvbert,placeholder-enstack,placeholder-dshgt}.

Reported accuracies on closed benchmarks routinely fall in the
$82$--$99\%$ range. However, several careful evaluations point out
that these scores degrade sharply when the test set is drawn from a
different project distribution \citep{placeholder-llm-detection-eval},
raising serious questions about the generalisation capability of
learning-based detectors beyond the corpora they were trained on.

A small number of studies use \emph{classical} machine learning rather
than deep learning, typically random forests or SVMs over hand-crafted
metrics, and report comparable performance on narrow targets while
remaining substantially cheaper to train. We mention them not because
they dominate the corpus, they do not, but because they often serve as
the ``simple baseline'' that the deep-learning contributions need to
beat.

\item \texttt{Hybrid Approaches}
\label{subsubsec:rw-hybrid}

A significant share of the corpus combines paradigms, and the works
split cleanly into two patterns. The first pattern chains a static
analyser with a learning-based component: GPTScan prompts an LLM with
patterns surfaced by static analysis on smart contracts
\citep{placeholder-gptscan}; PREV uses ML to triage SAST alerts;
MalCodeAI runs an autonomous agent that calls both static tooling and
an LLM for remediation \citep{placeholder-malcodeai}. The second
pattern fuses graph-based and sequence-based neural representations of
the same code: Anakin, GaGAT and GCRD report accuracies in the
$89$--$99\%$ range by feeding the model both an AST or PDG view and a
tokenised view of every function
\citep{placeholder-anakin,placeholder-gagat}.

These hybrid pipelines tend to dominate the high end of every reported
benchmark, but they also accumulate the limitations of their
components: brittleness to project-specific coding conventions on the
static side, opacity and dataset bias on the learning side. The
trade-off between accuracy and trustworthiness is, in our reading of
the corpus, the open question that distinguishes hybrid work from the
two purer paradigms above.

    
\end{enumerate}

\subsection{Automated Vulnerability Remediation}
\label{subsec:rw-apr}

The remediation literature is, by contrast, much more concentrated.
Automated Program Repair has been its dominant framing for the last
decade, and our mapping retains a comparatively small but
representative sample of five APR studies focused specifically on
security patches, complemented by the LLM-based repair entries already
discussed above. The discussion below mirrors the three-stage pipeline
introduced in Section~\ref{subsec:apr}: fault localization, patch
generation and patch validation.

\begin{enumerate}


 \item \texttt{Fault Localization (FL)}

Two broad regimes coexist. \emph{Traditional} FL collects the manual
techniques a developer applies, instrumenting the program with print
statements, stepping through it with a debugger. \emph{Advanced} FL,
the regime relevant to APR, can be grouped into six families. Slice-
and dice-based approaches narrow the search to the program slice
implicated by a failing variable; spectrum-based approaches exploit the
Executable Statement Hit Spectrum to compute a suspiciousness index per
statement; statistics-based approaches rank predicates by the
divergence of their truth distributions between passing and failing
runs. Program-state approaches inject snapshots from passing runs into
failing ones to test causality. Learning-based FL uses neural networks
(BP, RBF) or decision trees such as C4.5 to score statement
suspiciousness, while data-mining-based approaches extract $N$-grams or
association rules from execution traces, sometimes coupled with Formal
Concept Analysis. Model-based FL closes the list, comparing program
state against a reference model derived from manual rules or from
static analysis.

 \item \texttt{Patch Generation (PG)}

Four families dominate the PG literature. \emph{Search-based}
approaches explore the patch space heuristically, using genetic
programming on AST mutations or game-theoretic search guided by an
LTL specification. \emph{Constraint-based} approaches recast the
problem as constraint solving, leveraging contracts, symbolic
execution or separation logic to derive a constraint network whose
solutions are admissible patches. \emph{Learning-based} approaches
have grown rapidly: neural machine translation systems treat
bug-fixing as code-to-code translation, deep models work on the AST
rather than the textual form, and recent pipelines based on large
language models generate full candidate patches in a single call,
sometimes outperforming classical APR on standard benchmarks but
suffering from the same overfitting amplified by training-data
contamination \citep{placeholder-llm-repair}.
\emph{Pattern-based} approaches use either manually curated patch
templates (often mined from open-source projects) or automatically
extracted templates derived from change history; CNN-based pattern
extraction is one of the more recent variants in this family.

Within our security-focused sample, AVATAR and Sorald exploit fix
patterns of static-analysis violations to suggest patches at the
SonarQube level
\citep{placeholder-avatar,placeholder-sorald}; AddressWatcher localises
memory-leak fixes through sanitiser feedback
\citep{placeholder-addresswatcher}; PRD recompiles partial binaries to
mitigate vulnerabilities without source. The reported success rates
fall between $34\%$ and $65\%$ of correct or partially-correct
patches, with the wide variance directly tied to dataset composition
and to the strictness of the patch-correctness criterion used.

 \item \texttt{Patch Validation (PV)}

PV is the least mature of the three stages. Most works rely on the
project's own test suite as an oracle, which leaves the well-known
\emph{patch overfitting} problem open: a patch that passes the test
suite is not necessarily semantically correct, and most evaluations
acknowledge this caveat. Contract-based PV, mutation-based PV and
human-in-the-loop schemes have been proposed, but no single technique
is consensual at the time of writing.

   
\end{enumerate}
\subsection{Integrated Detection-Remediation Frameworks}
\label{subsec:integrated-rw}

Most works treat detection and remediation as two separate problems.
Detection tools emit warnings; remediation tools consume bug reports or
test failures. The number of frameworks that close the loop, that is,
that take an unannotated codebase as input and produce a corrected
codebase as output with no human in the middle, remains modest.
Sequencer \citep{placeholder-sequencer}, Hoppity
\citep{placeholder-hoppity} and a handful of recent LLM-driven pipelines
fall in this category for functional bugs; on the security side,
MalCodeAI \citep{placeholder-malcodeai} is one of the few systems that
genuinely chains autonomous detection and remediation, and most
``end-to-end'' systems still rely on a pre-existing detector from the
SAST family to feed the repair stage.

This sparsity is one of the gaps that motivates the present work. We
argue, and substantiate in the next chapter, that detection and
remediation are not merely two phases to be chained: the choice of
detector constrains the form of the patches that can be produced, and
the characteristics of the repair backend in turn influence which
detector outputs are actionable. Treating the two stages jointly is
therefore not just an engineering convenience but a research question
in its own right.

\textcolor{green}{correct all the references of this section}

\section{Synthesis \& Conclusion}
\label{sec:research-problem}

The overview presented in the previous sections brings several
limitations to the surface, and they affect the two halves of our
scope in different but compounding ways.

On the detection side, techniques have multiplied without a matching
effort on comparability. Static analysers, taint trackers and the
various AI-based detectors are evaluated on heterogeneous benchmarks,
often using incompatible definitions of vulnerability classes. Reported
accuracies above $95\%$ on internal datasets frequently degrade by
twenty points or more on cross-project evaluations, raising serious
questions about generalisation. The AI- and learning-based family
compounds this with opacity and dataset bias: many headline results
turn out to be partly explained by lexical artefacts of the training
data, and very few detectors expose a usable explanation of why a
given fragment was flagged.

On the remediation side, Automated Program Repair has matured
considerably, but its transposition from functional bug-fixing to
security-relevant flaws is far from straightforward. Patch overfitting
remains the rule rather than the exception: the standard test-based
oracle is a poor proxy for semantic correctness, and the semantic
preservation guarantees often advertised in the literature are weaker
than they appear once the patch is examined manually. LLM-based repair
compounds the problem with training-data contamination, where the test
bug or its fix is already present in the model's pre-training corpus.

The most striking gap, however, is structural and lies between the two
halves. Detection and remediation are studied as if they were
independent problems, even though an end-to-end secure-software
pipeline requires them to operate in concert: the form of a
detector's output constrains the form of the patches that can be
produced, and the expressiveness of the repair backend constrains
which detector outputs are actionable. The few works that bridge the
two stages do so for narrow vulnerability classes or on a single
language, and no consolidated map is available to guide the design of
more general pipelines.

These observations motivate the present work. The contribution
presented in the next chapter addresses the detection side of this gap
by proposing an approach for identifying software vulnerabilities in
source code, so that weaknesses can be surfaced and classified before
they propagate into production. The chapter that follows it then takes
up the remediation side, introducing an engine that repairs an
identified vulnerability only when the repair can be proven sound and
that abstains, with an explicit justification, otherwise. Taken
together, the two contributions are designed to be composed: the
detector fixes the form of the findings that the remediation stage
consumes, and the remediation stage closes the loop that the literature
surveyed above leaves open.